\chapter{MicroBooNE}

\section{The Short-Baseline Neutrino Program at Fermilab}

The Short-Baseline Neutrino (SBN) program at Fermilab places three liquid-argon
time projection chamber detectors on the Booster Neutrino Beam in a staged
near--middle--far arrangement. The point is not merely to instrument one
beamline several times, but to turn a common argon target and a common
reconstruction technology into an oscillation measurement in which many flux
and cross-section uncertainties are constrained \emph{in situ} rather than
imported from simulation alone. In that sense SBN is both a sterile-neutrino
search and a test of whether precision multi-detector LArTPC measurements can
be made with strongly correlated systematics~\cite{Brailsford:2017rxe,MicroBooNEProposal,Acciarri2017Future}.

Here the adjective \emph{short-baseline} is relative to the oscillation scale
rather than an absolute distance. Reactor disappearance searches historically
operated at source--detector separations of order
\(10\text{--}100\,\mathrm{m}\); for accelerator beams at hundreds of MeV to GeV
energies, comparable \(L/E\) values arise at baselines of a few hundred metres,
which is why the \(110\text{--}600\,\mathrm{m}\) SBN layout still sits
naturally in the short-baseline regime~\cite{GiuntiKim2007,Brailsford:2017rxe}.

This geometry matters because short-baseline oscillation signatures are subtle:
one is not looking for a gross change in topology, but for flavour- and
energy-dependent distortions that emerge only when the spectra measured at one
baseline are compared carefully with those at another.

The near detector fixes the incident beam and constrains interaction modelling before oscillation
effects can grow; the intermediate detector probes the region in which excesses
and mismodelling are most entangled; the far detector extends the lever arm so
that a genuine baseline-dependent effect cannot easily be mimicked by a single
mis-normalised flux or cross-section component. Because all three detectors use
liquid argon, the comparison is cleaner than in heterogeneous programs where
target-material differences must be unfolded alongside oscillation effects.

Even in the absence of new oscillation physics, SBN is a substantial detector
and interaction programme. The beam energies, detector granularity, and event
samples make it a powerful laboratory for neutrino--argon scattering, hadronic
final-state reconstruction, calorimetry, light collection, and reconstruction
validation in surface LArTPCs. The same measurements that stabilise the
oscillation comparison also sharpen the detector and interaction inputs needed
for future long-baseline argon experiments.

\input{figures/fermi-photo}

The three detectors play distinct but complementary roles:
\begin{itemize}
    \item \textbf{SBND} (Short-Baseline Near Detector) is the near detector,
    located \(110\,\mathrm{m}\) from the beam target. Its task is to measure
    the incident, effectively unoscillated beam directly in argon and to anchor
    the flux, topology, and interaction constraints that propagate downstream.
    In practice SBND is the control instrument of the program: it turns beam
    and cross-section uncertainties into measured quantities before they are
    reinterpreted as possible oscillation signals at longer baselines~\cite{Brailsford:2017rxe}.

    \item \textbf{MicroBooNE} is the intermediate detector, located
    \(470\,\mathrm{m}\) from the target. It occupies the most ambiguous
    position in the system: close enough that the beam remains intense and
    richly characterisable, but far enough that a genuine short-baseline effect
    could begin to deform the spectrum. That makes it a bridge between detector
    development and oscillation inference. It is also uniquely valuable as a
    dedicated imaging test of the low-energy electromagnetic excess first
    reported by MiniBooNE, because the LArTPC can separate electron-like and
    photon-like activity event by event rather than relying on Cherenkov-like
    ring classification alone~\cite{MicroBooNEProposal,MicroBooNEDesign,Acciarri2017Future}.

    \item \textbf{ICARUS} is the far detector, located \(600\,\mathrm{m}\)
    from the target. As the longest-baseline and largest-mass detector in SBN,
    it provides the strongest handle on flavour and spectral distortions once
    the near-detector constraints are propagated forward. It also brings into
    the program the mature large-volume LArTPC experience of the ICARUS line,
    linking the SBN effort to the broader evolution of argon-detector
    technology from demonstrator scale to precision oscillation instrument~\cite{Amerio2004,ICARUSPurity}.
\end{itemize}

\section{The MicroBooNE Detector}

\subsection{Liquid Argon Time Projection Chamber Technology}

Rubbia's 1977 argument for liquid argon as a neutrino target rested on a
conjunction of detector properties rather than on any single advantage. In
liquid form, argon is~\cite{Rubbia1977}:
\begin{enumerate}
\item it is \emph{dense}, with a liquid density of \SI{1.396}{\gram\per\cubic\centi\metre} at \qty{87}{\K}---it is over a third denser than water (\SI{1.0}{\gram\per\cubic\centi\metre}) and nearly three orders-of-magnitude denser than its own gaseous state (\SI{0.00178}{\gram\per\cubic\centi\metre})---enabling more compact detectors of comparable neutrino interaction probabilities~\cite{NISTLArDensity}.
\item it is \emph{electron-transparent}: when purified to oxygen-equivalent impurity concentrations below \qty{100}{ppt}, effective drift lengths exceeding \qty{5}{\m} have been achieved at a drift velocity of \qty{1.6}{\mm\per\micro\s} (\qty{500}{\V\per\cm}, \qty{87}{\K}); diffusion is modest, with longitudinal and transverse coefficients of \qty{4.8}{\centi\metre\squared\per\s} and \qty{13}{\centi\metre\squared\per\s}, yielding spreads of \qty{1.7}{\mm} and \qty{2.9}{\mm} over that distance---preserving millimetric spatial resolution.
\item it is \emph{comparatively inexpensive}: as a by-product of large-scale air-separation plants, bulk argon is available at \SIrange{140}{500}{\USD\per\tonne} depending on purity and supply, enabling kilotonne-scale detectors without prohibitive cost~\cite{ArgonPrice}.
\item it is \emph{easy to obtain and purify}: industrial air-separation units and standard molecular-sieve and getter filtration routinely deliver oxygen-equivalent impurity levels below \qty{100}{ppt}---over nine orders of magnitude lower than the \qty{21}{\%} oxygen concentration of ambient air---and suppress nitrogen contamination to the ppm level, more than six orders of magnitude beneath its \qty{78}{\%} atmospheric abundance, supporting multi-metre electron drift~\cite{ProtoDUNE2020,ICARUSPurity}.
\item it is \emph{chemically inert}: its closed-shell configuration ($3s^{2}3p^{6}$) yields a high ionisation energy and negligible electron affinity, providing no energetic incentive to form bonds; consequently it avoids corrosion or contamination of detector materials and cryogenic systems, minimising electronegative impurities that would otherwise limit electron lifetime~\cite{Rubbia1977}.
\end{enumerate}
Taken together, these properties motivate a detector in which the target,
tracking medium, and calorimeter are the same physical volume. The liquid-argon
time projection chamber is the mature realisation of that idea.

The original design already contains the essential modern picture: ionisation
electrons are drifted under a uniform electric field towards segmented sensor
planes, and the induced signals are combined into image-like projections of the
event. Subsequent developments have transformed that concept into a large,
stable cryogenic detector, but they have not changed its basic logic.

A liquid argon time projection chamber combines a dense neutrino target with
fine-grained electronic imaging. When a charged particle traverses the
detector, it both ionizes and excites the argon. The excitation produces
prompt scintillation light, while the ionization electrons drift under an
applied electric field toward an anode readout. If the interaction time
\(t_0\) is obtained from the light system and the drifting charge reaches the
anode at time \(t\), then the coordinate along the drift direction is
\[
x = v_d (t - t_0),
\]
where \(v_d\) is the electron drift velocity.

The remaining coordinates are obtained from the pattern of signals on multiple
readout planes, so the event can be reconstructed as a three-dimensional image
with calorimetric information attached.

\begin{figure}[htbp!]
    \centering
    \includegraphics[width=\textwidth]{figures/lartpc-concept.pdf}
    \caption{Conceptual schematic of a liquid argon time projection chamber, showing charge drift to segmented anode readout and complementary optical timing.}
    \label{fig:lartpc-concept}
\end{figure}

In MicroBooNE, the TPC consists of a cathode plane, a field cage that shapes
the drift field, and three wire planes on the anode side. The wire planes are
oriented at different angles, so each plane provides a complementary
projection of the same interaction. Combining those projections with the drift
time yields a full spatial reconstruction of vertices, tracks, and
electromagnetic showers.

This geometry is shown explicitly in
Fig.~\ref{fig:wire-plane-projection}. The collection plane defines the
\(W\) view, while the two induction planes define the \(U\) and \(V\) views
through wire orientations at \(\pm 60^\circ\). A localised charge deposition
therefore appears as three related two-dimensional footprints sharing a common
drift coordinate but different transverse coordinates, which is the geometric
basis of the pattern-recognition problem in MicroBooNE.

\input{figures/wire-plane-projection}

A separate optical system records the prompt scintillation light and fixes
\(t_0\), which is essential both for determining the drift coordinate and for
rejecting out-of-time cosmogenic activity in a surface detector.

The appeal of LArTPC technology for GeV-scale neutrino physics is that
tracking, vertexing, calorimetry, and particle identification are all
performed in the same homogeneous active volume. Fine spatial sampling gives
detailed event topology, and the collected charge provides calorimetric
information.

The local ionization density near the beginning of an electromagnetic shower
also allows electron--photon separation. That last point is central for
MicroBooNE, since the MiniBooNE low-energy excess could arise from either
electron-like or photon-like final states.

These advantages come with stringent detector requirements. The argon must be kept at cryogenic temperature and, more importantly, sufficiently pure that drifting electrons survive to the anode without excessive attachment to electronegative impurities. Stable high voltage, a uniform drift field, low-noise electronics, and efficient light collection are therefore part of the detector concept itself rather than auxiliary details. In practice, successful LArTPC operation is a coupled problem in cryogenics, instrumentation, calibration, and reconstruction.

The Micro Booster Neutrino Experiment (MicroBooNE) is an LArTPC located at Fermi National Accelerator Laboratory. It is part of the short-baseline neutrino (SBN) program and sits \(470\,\mathrm{m}\) downstream of the Booster Neutrino Beam (BNB) target, with an \(85\)-ton active volume instrumented by three wire planes and photomultiplier tubes~\cite{MicroBooNEDesign}.

\subsubsection{Detector Microphysics}

Consider the energy loss of a charged particle passing through liquid argon.
The total loss may be decomposed into electronic stopping, driven by inelastic
interactions with medium electrons, and nuclear stopping, driven by elastic
interactions with medium nuclei. In the regime most relevant to MicroBooNE
calorimetry, the visible detector response is governed predominantly by the
electronic component~\cite{ThomasImel1987,Doke2002}.

The electronic stopping power is expended through ionisation, excitation, and
the kinetic energy of the thermal electron. Writing the mean work required to
produce an electron--ion pair as \(W_{\mathrm{ion}}\), to avoid conflict with
the \(W\) wire view, one may express it as
\begin{equation}
  W_{\mathrm{ion}}
  =
  \bar E_i
  +
  \bar E_{\mathrm{ex}}\frac{\bar N_{\mathrm{ex}}}{\bar N_i}
  +
  \epsilon_i,
  \label{eq:ch3-work-function}
\end{equation}
where \(\bar E_i\) and \(\bar E_{\mathrm{ex}}\) are the mean energies expended
in ionisation and excitation, \(\bar N_i\) and \(\bar N_{\mathrm{ex}}\) are
the corresponding mean multiplicities, and \(\epsilon_i\) is the mean kinetic
energy of the thermal electron~\cite{Doke2002,Hitachi1983}. This relation sets
the initial electron--ion-pair yield,
\(N_e^0 \simeq E_{\mathrm{dep}}/W_{\mathrm{ion}}\), before recombination and
drift losses are applied.

The secondary electrons produced in inelastic collisions eventually thermalise
and may drift away from the parent ion before collection. That separation
affects electron--ion recombination and introduces a dependence on the local
linear energy transfer, commonly represented by the local \(dE/dx\). A
phenomenological parameterisation of the recombination probability is
\begin{equation}
  r
  =
  \frac{A\,dE/dx}{1 + B\,dE/dx}
  +
  C,
  \qquad
  C = 1 - \frac{A}{B},
  \label{eq:ch3-recombination-probability}
\end{equation}
where \(A\), \(B\), and \(C\) are recombination coefficients with an
electric-field dependence~\cite{Birks1964,ThomasImel1987}. Recombination is
stochastic, so it produces the familiar anti-correlation between the
ionisation and scintillation yields~\cite{Hitachi1983,Doke2002}. Charge
recombination limits the precision of calorimetric measurements. The charge
and light signals can be combined in a joint measurement that acts to reduce
this effect. In the MicroBooNE detector model used later in this chapter, that
generic recombination physics is implemented through the Modified Box
parametrisation summarised in Table~\ref{tab:fcl-allsettings}.

Scintillation itself carries a non-trivial time structure. The ionisation and
excitation channels populate the lowest singlet and triplet excimer states of
the self-trapped dimer \(\mathrm{Ar}_2^*\), which decay radiatively as
\begin{equation}
  \mathrm{Ar}_2^*
  \to
  2\,\mathrm{Ar} + \gamma,
  \label{eq:ch3-excimer-radiative-decay}
\end{equation}
producing vacuum-ultraviolet light near \(127\text{--}128\,\mathrm{nm}\) and
leaving the medium effectively transparent to its own scintillation
band~\cite{Hitachi1983,Doke2002}. The singlet component is prompt, while the
triplet component is delayed. For low-linear-energy-transfer depositions the
emission profile is therefore often approximated by
\begin{equation}
  I(t)
  =
  \frac{\alpha_s}{\tau_s}e^{-t/\tau_s}
  +
  \frac{\alpha_t}{\tau_t}e^{-t/\tau_t},
  \label{eq:ch3-scintillation-time-profile}
\end{equation}
where \(\alpha_s\) and \(\alpha_t\) are the singlet and triplet fractions and
\(\tau_s\) and \(\tau_t\) are the corresponding lifetimes. The late component
is suppressed when collisional quenching becomes efficient, so the
prompt-to-delayed ratio acquires a dependence on the local ionisation density.
That dependence underlies pulse-shape discrimination in liquid argon, even
though it is not the primary optical handle used in MicroBooNE event
reconstruction~\cite{Hitachi1983}.

More generally, xenon doping opens collisional transfer channels,
\begin{equation}
  \mathrm{Ar}_2^* + \mathrm{Xe}
  \to
  (\mathrm{ArXe})^* + \mathrm{Ar},
  \label{eq:ch3-arxe-transfer}
\end{equation}
\begin{equation}
  (\mathrm{ArXe})^* + \mathrm{Xe}
  \to
  \mathrm{Xe}_2^* + \mathrm{Ar},
  \label{eq:ch3-xe2-transfer}
\end{equation}
which compete with triplet quenching and reprocess part of the scintillation
into the xenon emission band near \(175\,\mathrm{nm}\)~\cite{McFadden2021}.
The resulting light is transported more favourably through the detector
because the longer wavelength suffers less Rayleigh scattering and is easier
to detect directly. Xenon-doped argon therefore exhibits a markedly different
time profile, with a reduced argon late tail and partial recovery of light
that would otherwise remain trapped in quenched triplet channels. MicroBooNE
does not operate with xenon-doped argon, but these reactions make clear that
the scintillation kinetics of LAr detectors can, in principle, be engineered.

\subsubsection{The Light Collection System}

The light collection system provides the prompt timing information that fixes
\(t_0\) for drift reconstruction and supports beam-coincidence and flash
matching in a surface detector. Liquid argon scintillation is emitted in the
vacuum-ultraviolet region around \(128\,\mathrm{nm}\), outside the direct
sensitivity band of the photomultiplier tubes used in MicroBooNE. Each primary
optical unit therefore places a tetraphenyl butadiene (TPB)-coated acrylic
plate in front of a \(200\,\mathrm{mm}\) PMT, so that the shifted visible
light can be detected with stable gain. A surrounding magnetic shield reduces
the effect of the ambient field on the tube response~\cite{MicroBooNEDesign}.

MicroBooNE deploys 32 of these optical units on the PMT rack mounted behind the
anode, distributed along the detector length. Four lightguide paddles were
also installed as a secondary R\&D system, but they are not used in the
analysis developed in this thesis~\cite{MicroBooNEDesign}. Together these
components provide the prompt-light information used for beam-gate timing,
flash matching, and the rejection of out-of-time cosmogenic activity.

Figure~\ref{fig:microboone-cryostat-schematic} places these subsystems in the
full cryostat geometry. The cutaway view shows the field-cage volume between
the APA and CPA sides of the TPC and the position of the PMT rack behind the
anode.

\input{figures/microboone-cryostat-schematic}

Figure~\ref{fig:microboone-optical-unit} then resolves a single optical unit
into the TPB-coated plate, photomultiplier tube, magnetic shield, and support
hardware that implement this light-collection channel.

\input{figures/microboone-optical-unit}

\subsubsection{The Space-Charge Effect}

For a near-surface LArTPC, the drift field is perturbed not only by the
electrode geometry but also by the persistent positive-ion population produced
by cosmogenic activity. These argon ions drift orders of magnitude more slowly
than the ionisation electrons, so they accumulate in the active volume and
distort the nominally uniform electric field. The result is a bias in the
time-to-position mapping and a corresponding deformation of reconstructed
tracks, with spatial offsets that reach the cm to \(10\) cm scale and approach
\(\sim 15\,\mathrm{cm}\) in the most affected regions of MicroBooNE~\cite{MicroBooNE_SCE_JINST2020}.

MicroBooNE derives correction maps from cosmic
muons~\cite{MicroBooNE_SCE_JINST2020}. The UV-laser system provides an
independent cross-check of the field model. These corrections are applied
during reconstruction, so later selections and kinematic observables are
defined in a space-charge-corrected coordinate system.

\begin{table}[htbp!]
    \centering
    \small
    \setlength{\tabcolsep}{6pt}
    \caption{Key MicroBooNE experiment parameters (site, beams, and detector design).}
    \label{tab:uboone_overview}
    \begin{tabular}{@{}p{0.38\linewidth}p{0.58\linewidth}@{}}
        \toprule
        \textbf{Quantity} & \textbf{Value} \\
        \midrule
        \multicolumn{2}{@{}l@{}}{\textbf{Site and beams}} \\
        Location & Fermilab (Batavia, IL), Liquid Argon Test Facility (LArTF) \\
        BNB baseline & \(470\,\mathrm{m}\) downstream of the BNB target (on-axis) \\
        BNB spectrum & Predominantly \(\nu_\mu\) from \(\pi\) decays; energy peak \(\sim 700\,\mathrm{MeV}\) \\
        NuMI exposure & Off-axis component; \(\langle E_\nu\rangle \sim 0.25\,\mathrm{GeV}\) (\(\pi\)) and \(\sim 2\,\mathrm{GeV}\) (\(K\)); \(\sim 600\,\mathrm{m}\) downstream of NuMI target \\

        \addlinespace[0.4em]
        \multicolumn{2}{@{}l@{}}{\textbf{LArTPC}} \\
        LArTPC dimensions (V\(\times\)H\(\times\)L) & \(2.325\,\mathrm{m}\times 2.560\,\mathrm{m}\times 10.368\,\mathrm{m}\) \\
        LArTPC argon mass & \(90\,\mathrm{t}\) \\
        Drift field & \(500\,\mathrm{V/cm}\) \\
        Drift high voltage & Cathode at \(-128\,\mathrm{kV}\) (nominal) \\
        Anode plane bias voltage & \(-110\,\mathrm{V},\,0\,\mathrm{V},\,+230\,\mathrm{V}\) (U,V,Y) \\
        Max drift length & \(2.560\,\mathrm{m}\) \\
        Drift velocity & \(1.14\,\mathrm{mm}/\mu\mathrm{s}\) \\
        Max drift time & \(1.6\,\mathrm{ms}\) (cathode \(\rightarrow\) first induction plane, at \(500\,\mathrm{V/cm}\)) \\
        Readout planes & 3 wire planes (U,V induction; Y collection) \\
        Wire pitch / plane spacing & \(3\,\mathrm{mm}\) / \(3\,\mathrm{mm}\) \\
        Wire angles (U,V,Y) & \((+60^\circ,\,-60^\circ,\,0^\circ)\) w.r.t.\ vertical \\
        Wires & 8256 total (U: 2400; V: 2400; Y: 3456) \\
        Wire type & Stainless steel \\
        Wire diameter & \(150\,\mu\mathrm{m}\) \\
        Wire coating & \(2\,\mu\mathrm{m}\) Cu, \(0.1\,\mu\mathrm{m}\) Ag \\

        \addlinespace[0.4em]
        \multicolumn{2}{@{}l@{}}{\textbf{Cryostat and operating conditions}} \\
        Total liquid argon mass & \(170\,\mathrm{t}\) \\
        Operating temperature / pressure & \(87\,\mathrm{K}\) / \(1.24\,\mathrm{bar}\) \\
        Cryostat geometry & \(12.2\,\mathrm{m}\) length; \(3.81\,\mathrm{m}\) inner diameter \\

        \addlinespace[0.4em]
        \multicolumn{2}{@{}l@{}}{\textbf{Light collection}} \\
        Primary system & 32 optical units with 200\,mm (8\,in) PMTs (0.9\% photocathode coverage) \\
        Secondary system & 4 lightguide paddles (R\&D) \\
        \bottomrule
    \end{tabular}
\end{table}

Table~\ref{tab:fcl-allsettings} summarises the key MicroBooNE detector-property, LAr material, and simulation settings used throughout this analysis. These reflect the standard MicroBooNE detector-services configurations in \texttt{uboonecode}.

\begin{table}[htbp!]
    \centering
    \small
    \setlength{\tabcolsep}{5pt}
    \caption{Consolidated MicroBooNE detector-property, LAr material, and simulation settings used in this analysis.}
    \label{tab:fcl-allsettings}
    \begin{tabular}{@{}l l l p{0.48\linewidth}@{}}
        \toprule
        Source & Parameter & Value & Comment / units \\
        \midrule
        \multicolumn{4}{@{}l}{\textbf{DetectorPropertiesServiceStandard (\texttt{standard\_detproperties})}} \\
        \multicolumn{4}{@{}l}{\emph{Drift and charge transport}} \\
        detprop & \texttt{Temperature}      & \(91\,\mathrm{K}\) & LAr temperature. \\
        detprop & \texttt{Efield}           & 0.50, 0.66, 0.80 & Electric field(s) in \(\mathrm{kV/cm}\). \\
        detprop & \texttt{Electronlifetime} & \(750\,\mu\mathrm{s}\) & Electron lifetime. \\

        \addlinespace
        \multicolumn{4}{@{}l}{\emph{Electronics and readout window}} \\
        detprop & \texttt{ElectronsToADC}    & \(1.208041\times 10^{-3}\) & ADC/e\(^{-}\) (as configured). \\
        detprop & \texttt{NumberTimeSamples} & 2048 & Ticks per event readout frame. \\
        detprop & \texttt{ReadOutWindowSize} & 2048 & Ticks in readout window. \\
        detprop & \texttt{TimeOffsetU}       & -5.193 & Tick offset (U view/plane). \\
        detprop & \texttt{TimeOffsetV}       & 0.585  & Tick offset (V view/plane). \\
        detprop & \texttt{TimeOffsetZ}       & 0      & Tick offset (third view/plane). \\

        \addlinespace
        \multicolumn{4}{@{}l}{\textbf{LArPropertiesServiceStandard (\texttt{standard\_properties})}} \\
        \multicolumn{4}{@{}l}{\emph{Material constants}} \\
        larprop & \texttt{RadiationLength}  & \(19.55\,\mathrm{g/cm^2}\) & Radiation length (mass thickness). \\
        larprop & \texttt{AtomicNumber}     & 18 & Argon \(Z\). \\
        larprop & \texttt{AtomicMass}       & \(39.948\,\mathrm{g/mol}\) & Argon molar mass. \\
        larprop & \texttt{ExcitationEnergy} & \(188.0\,\mathrm{eV}\) & Mean excitation energy. \\

        \addlinespace
        \multicolumn{4}{@{}l}{\textbf{Geant4/LArTPC simulation (\texttt{standard\_largeantparameters})}} \\
        \multicolumn{4}{@{}l}{\emph{Charge transport (diffusion)}} \\
        larG4 & \texttt{LongitudinalDiffusion} & \(6.2\times 10^{-9}\)  & \(\mathrm{cm^2/ns}\). \\
        larG4 & \texttt{TransverseDiffusion}   & \(16.3\times 10^{-9}\) & \(\mathrm{cm^2/ns}\). \\

        \addlinespace
        \multicolumn{4}{@{}l}{\emph{Recombination model (parameters)}} \\
        larG4 & \texttt{ModBoxA} & 0.930 & Modified Box \(A\) (alpha). \\
        larG4 & \texttt{ModBoxB} & 0.212 & Modified Box \(B\) (beta; field-scaled in implementation). \\
        \bottomrule
    \end{tabular}
\end{table}

\subsection{Data and Simulation Samples}

An event in this analysis is a single triggered detector readout: the full
wire-waveform window together with the coincident optical record. In offline
processing this record is represented by one \texttt{art::Event} identified by
its \texttt{(run, subrun, event)} tuple, so it is not synonymous with a single
neutrino interaction. That distinction matters in MicroBooNE because both the
data and the simulation are reconstructed as complete detector readouts rather
than as isolated beam interactions.

Figure~\ref{fig:microboone-readout-timing} shows the timing structure that
defines this event record. The TPC and PMT systems are read out on
millisecond-scale windows around the beam trigger, while the all-PMT beam gate
that tags beam-coincident optical activity occupies only a
\SI{23.4}{\micro\second} interval inside that larger record. The separation of
these timescales is one of the reasons cosmic activity remains an intrinsic
part of the MicroBooNE reconstruction problem.

\input{figures/microboone-readout-timing}

The timing structure specifies when the detector is sampled; the readout
architecture specifies how those signals are carried to the data-acquisition
system. Figure~\ref{fig:microboone-readout-scheme} shows the separation between
the TPC and optical chains, from in-cryostat front-end electronics through the
warm interface and detector-hall digitisation, and makes explicit where the
beam-gate trigger enters the PMT path.

\begin{figure}[htbp!]
    \centering
    \includegraphics[width=\textwidth]{figures/microboone-readout-scheme-cropped.pdf}
    \caption{MicroBooNE TPC and PMT readout scheme. The figure answers how wire and optical signals are routed from the cryostat to the detector-hall data-acquisition system: the TPC and PMT chains are processed separately, digitised in dedicated front-end modules, buffered, and passed onward through optical links, while the beam gate enters through the trigger logic on the optical side.}
    \label{fig:microboone-readout-scheme}
\end{figure}

The Monte Carlo samples used here follow the MCC9 overlay strategy. A simulated
beam neutrino interaction is propagated through the full detector simulation and
then embedded, channel by channel, into a real external-trigger cosmic readout.
If \(d_c(t)\) denotes the recorded EXT waveform on channel \(c\), \(m_c(t)\) the
simulated beam waveform, \(p_c^{\mathrm{MC}}\) the simulated pedestal, and
\(g_c\) an optional channel gain factor, the overlaid waveform is
\begin{equation}
  o_c(t)
  =
  d_c(t)
  +
  g_c\!\left[m_c(t)-p_c^{\mathrm{MC}}\right].
  \label{eq:ch3-overlay-waveform}
\end{equation}
The resulting event is then reconstructed with the same chain used for data.
Overlay therefore preserves the data noise environment and the real cosmic
activity, while still providing a truth-tagged beam component for efficiency
and background studies.

The sample classes entering the analysis are:
\begin{description}
  \item[Beam-on overlay.] Inclusive simulated \(\nu\)-argon interactions overlaid
  on external-trigger cosmic data and reconstructed with the standard data
  chain.
  \item[Strangeness-enriched overlay.] The same overlay construction, but with a
  generator-level filter that retains only interactions with at least one
  strange hadron at nuclear exit.
  \item[Dirt.] Simulated neutrino interactions outside the active TPC whose
  secondaries enter the detector volume, again overlaid on external-trigger
  data.
  \item[External.] Beam-off data taken with the same detector gate, used both as
  the cosmic template in the analysis and as the cosmic component of the overlay
  productions.
\end{description}

Because the cosmic component of an overlaid event has no corresponding truth
particle, truth matching is intrinsically partial. For any reconstructed object
\(T\) matched to a candidate truth particle \(p^\star\), the analysis therefore
requires a minimum backtracking purity,
\begin{equation}
  \mathrm{purity}(T)
  =
  \frac{Q_{p^\star}(T)}{Q(T)},
  \qquad
  \mathrm{accept}\iff \mathrm{purity}(T)>0.1,
  \label{eq:ch3-overlay-purity}
\end{equation}
where \(Q(T)\) is the total attributed charge of the object and \(Q_{p^\star}(T)\)
the subset assigned to \(p^\star\). This convention prevents the external cosmic
component from being misinterpreted as a valid truth match to the beam event.

The analysis uses the five MicroBooNE accelerator run periods from
1 October 2015 to 31 March 2020. These periods define the delivered exposures
used later for template normalisation for the BNB sample and for the NuMI
forward- and reverse-horn-current configurations. For the off-axis NuMI
measurement, the full dataset corresponds to \(\num{1.028e21}\) POT in forward
horn current and \(\num{1.210e21}\) POT in reverse horn current.

Template normalisation follows the delivered exposure. For beam samples the
scaling factor is
\begin{equation}
  S_{\mathrm{POT}}
  =
  \frac{\mathrm{POT}_{\mathrm{run}}}{\mathrm{POT}_{\mathrm{sample}}},
  \label{eq:ch3-pot-scale}
\end{equation}
while the beam-off external samples are normalised by the accepted trigger
count,
\begin{equation}
  S_{\mathrm{trig}}
  =
  \frac{N_{\mathrm{trig}}^{\mathrm{run}}}
       {N_{\mathrm{trig}}^{\mathrm{sample}}}.
  \label{eq:ch3-trigger-scale}
\end{equation}
These factors are absorbed into the nominal per-event weights used later in the
profile-likelihood fit.

The nominal beam predictions are generated with the standard MCC9
\texttt{GENIE}-based overlay chain, while the strangeness-enriched samples reuse
the same detector simulation and reconstruction but apply a generator-level
filter. The analysis-level orthogonality used to combine those enriched samples
with the inclusive Beam-On prediction is part of the measurement definition and
is discussed in Section~\ref{subsec:analysis-sample-orthogonality}. At detector
level, the relevant point here is simply that the enriched production preserves
the same response model while improving Monte Carlo support for rare
strangeness-bearing topologies.

\subsection{Wire Processing and Detector Images}
The MicroBooNE TPC does not present reconstruction with three-dimensional
points directly. It presents digitised waveforms on \num{8256} wires over a
\num{2048}-tick readout window, and the first reconstruction problem is to
compress that raw electronic record into a stable charge representation. In the
\texttt{LArSoft}/\texttt{art} chain used by MicroBooNE, this is done by noise
filtering, response deconvolution, region-of-interest finding, and hit
reconstruction~\cite{SniderPetrillo:2017LArSoft,ubnoise,ubsp1,ubsp2}. At
waveform level the measured signal on a given wire can be written
\begin{equation}
  M(t_0)=\int R(t-t_0)\,S(t)\,dt + n(t_0),
  \label{eq:ch3-waveform-convolution}
\end{equation}
as the convolution of the drifted-charge distribution \(S(t)\) with the
combined field-and-electronics response \(R(t)\), plus noise \(n(t_0)\). In
frequency space, deconvolution therefore acts as a regularised inversion,
\begin{equation}
  \hat{S}(\omega)=\frac{M(\omega)}{R(\omega)}\,F(\omega),
  \label{eq:ch3-fourier-deconvolution}
\end{equation}
where \(F(\omega)\) suppresses modes that would otherwise be dominated by noise
amplification.

The plane dependence of the response is shown in
Fig.~\ref{fig:wire-plane-waveforms}. The \(U\) and \(V\) induction wires
produce bipolar signals because the induced current changes sign as the charge
packet approaches and then passes the plane, whereas the \(W\) collection wire
records a predominantly unipolar pulse. That distinction is not merely
diagnostic: it is the reason induction-plane charge extraction is more
sensitive to non-local response modelling and low-frequency noise than the
collection-plane case~\cite{ubsp1,ubsp2,ubnoise}.

\input{figures/wire-plane-waveforms}

With sampled waveforms, Eq.~\eqref{eq:ch3-waveform-convolution} becomes the
linear system
\begin{equation}
  M_i=\sum_j R_{ij}S_j,
  \label{eq:ch3-discrete-convolution}
\end{equation}
and the stabilised solution may be viewed as a regularised least-squares
problem,
\begin{equation}
  \chi^2
  =
  \sum_i\left(M_i-\sum_j R_{ij}S_j\right)^2
  + c^2\sum_i (S_i'')^2,
  \qquad
  S_i''\simeq S_{i+1}-2S_i+S_{i-1},
  \label{eq:ch3-regularised-inversion}
\end{equation}
which penalises the high-frequency fluctuations most vulnerable to noise.

This step is especially delicate on the induction planes. Their bipolar
weighting response is extended in both time and neighbouring-wire index, so the
signal on one wire is not purely local to that wire. Schematically,
\begin{equation}
  M_i(\omega)=R_0(\omega)S_i(\omega)+R_1(\omega)S_{i+1}(\omega)+\cdots,
  \label{eq:ch3-2d-deconvolution}
\end{equation}
which motivates a joint inversion in wire and time for the induction planes,
whereas the collection-plane response is comparatively local in wire. The
bipolar induction response further suppresses the low-frequency component,
\begin{equation}
  R_{\mathrm{ind}}(\omega\to 0)\approx 0,
  \label{eq:ch3-induction-lowfreq}
\end{equation}
so MicroBooNE combines coherent-noise filtering with region-of-interest (ROI)
selection and adaptive-baseline correction before hit extraction. A standard
reference is the Wiener filter,
\begin{equation}
  F_{\mathrm{Wiener}}(\omega)=\frac{S^2(\omega)}{S^2(\omega)+N^2(\omega)},
  \label{eq:ch3-wiener}
\end{equation}
but in practice MicroBooNE uses a modified high-frequency roll-off filter for
charge extraction together with an additional low-frequency suppression filter
for ROI finding~\cite{ubsp1,ubsp2,ubnoise}. The output of that stage is a
calibrated collection of \texttt{recob::Hit} objects, each parameterised by
wire, peak time, width, and integrated charge. These hits are the practical
interface between detector readout and higher-level reconstruction: slicing,
flash matching, and Pandora all begin from this compressed representation.

The qualitative effect of that processing is visible in representative
MicroBooNE signal-processing displays. Figure~\ref{fig:deconv-event-examples}
shows a \(U\)-plane induction-view example, where the field response is most
non-local and the gain from deconvolution is correspondingly clearest. The
processed image sharpens near-vertical charge structure, makes short tracks and
delta rays near the interaction vertex more distinct, and suppresses diffuse
artefacts that would otherwise obscure the topology.

\begin{figure}[htbp!]
  \centering
  \includegraphics[width=0.94\textwidth]{figures/microboone-deconv-event-3493-u-cropped.pdf}
  \caption[Representative MicroBooNE induction-plane deconvolution display.]{Representative MicroBooNE induction-plane deconvolution display, chosen to answer what image-level structure is recovered where signal processing matters most. The internal subpanels compare the raw waveform image, the software noise-filtered image, and the final two-dimensional deconvolved image for a \(U\)-plane view. The progression shows that induction-plane deconvolution sharpens charge localisation near the neutrino vertex and improves the visibility of short prongs and delta-ray structure.}
  \label{fig:deconv-event-examples}
\end{figure}

A representative processed wire--time view is shown in
Fig.~\ref{fig:processed-event-display-bnb}. This display complements the
induction-plane deconvolution example by showing the sparse, topologically
structured image that later reconstruction must interpret.

\begin{figure}[htbp!]
  \centering
  \includegraphics[width=0.92\textwidth]{figures/microboone-event-display-run3493-collection.pdf}
  \caption[Representative MicroBooNE collection-plane event display.]{Representative MicroBooNE collection-plane event display after signal processing. The image shows the sparse wire--time structure presented to downstream reconstruction in a typical neutrino candidate, including compact vertex activity, extended track segments, and lower-level secondary structure.}
  \label{fig:processed-event-display-bnb}
\end{figure}

For the image-level analysis developed later in this thesis, the calibrated hit
record is not reduced immediately to a small set of reconstructed particles.
Instead, the selected neutrino slice is re-expressed as three fixed-size
wire--time images, one for each TPC view. The image window is centred by a
charge-weighted three-dimensional centroid \(\vec c_{\mathrm{world}}\) built
from slice-associated space points, with hit charge distributed across the
space points attached to each hit. To reduce the leverage of sparse outliers,
the final centre is taken from a trimmed core of the charge-weighted space-point
cloud rather than from the full event extent.

For each view \(\alpha\in\{U,V,Y\}\), this centre is projected into a
wire--drift coordinate system. Writing \(\theta_U=+60^\circ\),
\(\theta_V=-60^\circ\), and \(\theta_Y=0^\circ\), the local wire coordinate is
\begin{equation}
  w_\alpha = Z\cos\theta_\alpha - Y\sin\theta_\alpha,
\end{equation}
while the vertical image axis is the drift coordinate \(x\). With image
dimensions \((W,H)\), wire pitch \(\Delta w_\alpha\), and drift-step
\(\Delta x_{\mathrm{drift}}=v_{\mathrm{drift}}\Delta t_{\mathrm{tick}}\), the
per-view image origin is set by
\begin{equation}
  w_0 = w_c - \frac{W\Delta w_\alpha}{2},
  \qquad
  x_0 = x_c - \frac{H\Delta x_{\mathrm{drift}}}{2},
\end{equation}
so that detector coordinates are mapped to integer pixel indices by
\begin{equation}
  c(w)=\left\lfloor\frac{w-w_0}{\Delta w_\alpha}\right\rfloor,
  \qquad
  r(x)=\left\lfloor\frac{x-x_0}{\Delta x_{\mathrm{drift}}}\right\rfloor.
\end{equation}

The detector image itself is filled by looping over waveform regions of
interest associated with the selected hits. For each contributing ADC sample
\(a(t)\) on a wire in view \(\alpha\), the algorithm converts the tick to drift
coordinate, maps it to a row \(r(t)\), and accumulates the sample into the
corresponding pixel:
\begin{equation}
  I_{\mathrm{det}}^{(\alpha)}(r(t),c)
  \leftarrow
  I_{\mathrm{det}}^{(\alpha)}(r(t),c)
  + a(t)\,\mathbf{1}\!\left[a(t)>a_{\mathrm{thr}}\right].
\end{equation}
If multiple hits or multiple ROIs map to the same pixel, their charge
contributions add. This preserves a detector-native representation of
branching, displacement, local charge concentration, and partial overlap,
rather than forcing the event immediately into a particle-flow summary.

\subsubsection{Unresponsive Channels and Image Masks}

Not every MicroBooNE readout channel contributes a usable waveform in the
run-period configuration used for this analysis. The image chain therefore
propagates an explicit bad-channel mask and removes hits on masked channels
before either detector or semantic images are written. For the configuration
used here, the mask excludes 311 of the 8256 TPC channels: 63 in the \(U\)
plane, 51 in the \(V\) plane, and 197 in the \(Y\) plane.

Figure~\ref{fig:microboone-bad-channel-mask} summarises that mask in local
channel coordinates. Most of the induction-plane exclusions occur in short,
isolated runs, whereas the collection plane contains a broad inactive span.
Because channel status is a wire property rather than a drift-time property,
each masked channel becomes a blind column in the corresponding wire--time
image. The extended collection-plane gap is the dominant feature of the mask
and is one reason the later fiducial definition removes an inactive
longitudinal span from the signal volume
(Section~\ref{sec:analysis-signal-definition}).

\input{figures/microboone-bad-channel-mask}

In simulation, an accompanying semantic image is built on the same lattice. A
backtracker association identifies the dominant contributing MC particle for
each hit, assigns it a discrete semantic label \(\ell(h)\), and writes that
label to the corresponding pixel,
\begin{equation}
  I_{\mathrm{sem}}^{(\alpha)}(r,c)\leftarrow \ell(h),
\end{equation}
providing a view-aligned truth interpretation of the same detector image. The
detector and semantic images therefore share geometry and indexing, but differ
in what they encode: one carries measured charge, the other simulation truth.

\subsubsection{Drift Coordinate and Lifetime Correction}

The mapping from TPC ticks to image coordinates and the correction for
drift-time-dependent charge loss are handled uniformly in the image chain. For a
given plane, the nominal drift coordinate associated with readout tick
\(t_{\mathrm{tick}}\) is obtained from the detector-properties service as
\begin{equation}
  x_{\mathrm{nom}}(t_{\mathrm{tick}};\mathrm{plane})
  =
  \texttt{DetectorProperties::ConvertTicksToX}(t_{\mathrm{tick}},\mathrm{plane}),
  \label{eq:ch3-convert-ticks-to-x}
\end{equation}
which folds in the electronics sampling period, the drift velocity, and the
plane-dependent clock offsets. For a hit on a wire centred at
\((x_{\mathrm{wire}},y_{\mathrm{wire}},z_{\mathrm{wire}})\), the corresponding
nominal three-dimensional point is
\begin{equation}
  \mathbf p_{\mathrm{nom}}(t_{\mathrm{tick}})
  =
  \bigl(x_{\mathrm{nom}}(t_{\mathrm{tick}}),y_{\mathrm{wire}},z_{\mathrm{wire}}\bigr).
\end{equation}
After the space-charge and geometry corrections applied by the image utilities,
this point is projected into the chosen view and mapped into image coordinates.
The column index is set by the projected transverse coordinate, while the row
index for each waveform sample is determined by the corrected drift position,
\begin{equation}
  \mathrm{row}(t)
  =
  \texttt{ImageProperties::row}\bigl(x_{\mathrm{corr}}(t)\bigr).
  \label{eq:ch3-image-row}
\end{equation}

The event start time \(T_0\) enters through the calorimetric correction rather
than through the image geometry. For hit peak time \(t_{\mathrm{peak}}\) and
tick period \(\Delta t_{\mathrm{tick}}\), the approximate drift time is
\begin{equation}
  t_{\mathrm{drift}}
  \simeq
  t_{\mathrm{peak}}\,\Delta t_{\mathrm{tick}} - T_0.
  \label{eq:ch3-drift-time}
\end{equation}
This drift time is passed to the calorimetry algorithm so that the inferred hit
energy is corrected for electron-lifetime attenuation along the drift path. In
the notation of the reconstruction utilities,
\begin{equation}
  \left(\frac{dE}{dx}\right)_{\mathrm{hit}}
  =
  \texttt{CalorimetryAlg::dEdx\_AREA}(h,\mathrm{pitch},T_0),
  \qquad
  E_{\mathrm{hit}}
  =
  \left(\frac{dE}{dx}\right)_{\mathrm{hit}}\times \mathrm{pitch}.
  \label{eq:ch3-hit-energy}
\end{equation}
The energy \(E_{\mathrm{hit}}\) is then distributed over the image pixels
according to the local waveform amplitudes, so the final detector images encode
not raw ADC values alone but a lifetime-corrected, approximately calorimetric
charge density along the drift direction.

\subsection{Inference Deployment}
The detector images are not consumed inside the validated reconstruction stack
by linking the sparse-CNN code directly against the experiment release.
Instead, inference is deliberately decoupled from reconstruction. In the
MicroBooNE software environment the \texttt{LArSoft}/UPS release fixes the
compiler and library interfaces used by the process, while the PyTorch and
MinkowskiEngine stack used by the classifier evolves on a different cadence.
Treating inference as an external, containerised service therefore avoids
binding the physics release to a rapidly changing Python and GPU software
environment.

Operationally, the \texttt{ImageProduction} stage writes a minimal three-plane
request for each selected event in channel-first layout,
\[
x \in \mathbb{R}^{3\times H\times W},
\]
using the \(U\), \(V\), and \(Y\) detector images defined above. The
\texttt{InferenceProducer} module consumes these per-view images, identifies the
induction and collection planes through the encoded \texttt{geo::View\_t}
information, checks that the declared \((H,W)\) dimensions are consistent with
the flattened ADC buffers, and assembles the ordered image tuple
\((x_U,x_V,x_Y)\). Events whose three views are identically empty are
short-circuited at this stage and record no classifier output.

For non-empty requests, the producer writes the three planes to a binary file,
\texttt{planes.bin}, in single-precision channel-first layout. The scratch
directory used for this handoff is resolved at runtime: a configured
\texttt{ScratchDir} parameter takes precedence; otherwise the module falls back
to the batch scratch area, \texttt{\_CONDOR\_SCRATCH\_DIR}, if it is defined,
and finally to the current working directory. That directory is then
bind-mounted into an external Apptainer runtime together with the network
assets. In production the container image resides on the experiment CVMFS area,
specifically the \texttt{lantern\_v2\_me\_06\_03\_prod} image, and is invoked
with \texttt{apptainer exec --cleanenv}. Inside that environment a lightweight
Python wrapper reads \texttt{planes.bin}, performs the sparse-tensor
quantisation, constructs the \((\mathrm{view},y,x)\) coordinate set, runs the
trained network, and writes back both the classifier response and a compact
metadata record.

The reconstruction process then reads that response back into stable event
products. For each configured model, the output is parsed into a
\texttt{ModelPred} object holding the returned class scores and a
\texttt{ModelPerf} object holding the timing breakdown. The latter separates the
time to materialise the request (\(t_{\mathrm{write}}\)), the container
execution time (\(t_{\mathrm{exec}}\)), the response-parse time
(\(t_{\mathrm{read}}\)), and the corresponding setup, inference, and
post-processing timings reported from inside the container, together with the
peak resident set size. The physics analysis therefore sees only stable,
versioned prediction and performance products, while the implementation details
of the sparse-CNN stack remain isolated from the detector-simulation release.
This design keeps the image-level analysis reproducible at production scale and
allows the inference software to be updated, benchmarked, and redeployed
without perturbing the validated reconstruction chain.

\section{The NuMI Beam}

MicroBooNE also observes neutrinos from the Neutrinos at the Main Injector
(NuMI) beamline from an off-axis location. The centre of the TPC is
\(0.679\,\mathrm{km}\) downstream of the NuMI target and
\(91.1\,\mathrm{m}\) transversely displaced from the beam axis
(\(55.0\,\mathrm{m}\) horizontally and \(72.6\,\mathrm{m}\) vertically),
corresponding to an off-axis angle of \(7.71^\circ\).

That geometry is not a
small correction: it is one of the main ingredients that determines the energy
spectrum and flavour content relevant for the NuMI-based analysis developed
later in this thesis~\cite{Abramov:2002NuMIOptics,Adamson:2016NuMIBeam,Aliaga:2016NuMIPRD}.
Figure~\ref{fig:fermi-map} identifies the Booster and Main Injector beamlines
within the wider Fermilab accelerator complex.

\input{figures/fermi-map}

At the level of the accelerator site this establishes where the NuMI beam is
generated. The local beamline elements that shape the secondary hadrons and the
resulting neutrino trajectories are described next.

NuMI is a conventional horn-focused neutrino beam. A \(120\,\mathrm{GeV}\)
proton beam from the Main Injector strikes a segmented graphite target,
producing secondary pions and kaons that are sign-selected by a pair of
magnetic horns and transported into a \(675\,\mathrm{m}\) helium-filled decay
pipe.

The surviving hadrons are then absorbed, while the neutrinos produced in
meson and muon decays continue toward downstream detectors. In forward-horn
current running the horns preferentially focus positively charged mesons and
therefore enhance the \(\nu_\mu\) component; in reverse-horn current running
the polarity is reversed and the beam becomes \(\bar\nu_\mu\)-enhanced, albeit
with non-negligible wrong-sign contamination because the sign selection is
never perfect~\cite{Adamson:2016NuMIBeam}.

To leading order, the muon-flavour flux is generated by the two-body decays
\(\pi^\pm\to\mu^\pm\nu_\mu(\bar\nu_\mu)\) and
\(K^\pm\to\mu^\pm\nu_\mu(\bar\nu_\mu)\). For a parent meson
\(M\in\{\pi,K\}\) of mass \(m_M\), energy \(E_M\), momentum \(p_M\), and
laboratory decay angle \(\theta_\nu\) between the neutrino and the parent
direction, the neutrino energy is
\begin{equation}
  E_\nu(\theta_\nu;E_M)
  = \frac{m_M^2-m_\mu^2}{2\left(E_M-p_M\cos\theta_\nu\right)}
  \;\simeq\;
  \left(1-\frac{m_\mu^2}{m_M^2}\right)\frac{E_M}{1+\gamma_M^2\theta_\nu^2},
  \label{eq:numi-two-body-energy}
\end{equation}
where \(\gamma_M\equiv E_M/m_M\). This relation makes the off-axis geometry
physically useful: at fixed viewing angle the factor
\(1+\gamma_M^2\theta_\nu^2\) suppresses neutrinos from the highest-energy
parents, so MicroBooNE sees a comparatively narrower few-GeV spectrum than an
on-axis NuMI detector. The resulting off-axis beam still carries appreciable
contributions from both pion- and kaon-parent decays, which is one reason the
NuMI flux at MicroBooNE remains rich in both rate and spectral structure
despite the geometric suppression~\cite{Aliaga:2016NuMIPRD}.

Figure~\ref{fig:numi-horn-focusing} summarises the local target--horn--decay
geometry behind that picture, including representative focused and unfocused
parent trajectories and the corresponding neutrino lines of flight.

\input{figures/numi-horn-focusing}

For this thesis, the important practical point is that the NuMI prediction used
for signal and background modelling is not just a generic Fermilab beam
description but a MicroBooNE-specific off-axis flux model. The central
prediction is built from a \texttt{Geant4}-based beamline simulation combined
with \texttt{PPFX} hadron-production reweighting, and the updated MicroBooNE
configuration restores shielding material in the target hall, retunes the
hadronic transport, and propagates revised hadron-production uncertainties for
the off-axis acceptance.

These changes are not cosmetic. They produce energy-dependent modifications of
up to roughly \(40\%\) in the predicted
MicroBooNE off-axis flux in some regions of phase space, while remaining
consistent with the constrained on-axis NuMI predictions used by other
experiments~\cite{Aliaga:2016NuMIPRD,MicroBooNE:2025NuMIFlux}. The beam model
therefore enters the present analysis as a genuine physics input rather than as
an interchangeable normalisation constant.

\section{Pandora Reconstruction}

MicroBooNE reconstruction is organised around the multi-algorithm Pandora
framework~\cite{MarshallThomson:2015PandoraSDK,Acciarri:2018PandoraMicroBooNE}.
The point of Pandora is not to impose a single global fit on the event, but to
build and compare many local topological hypotheses until a consistent
particle-flow description emerges. That strategy is well suited to LArTPC data,
where pattern recognition is driven as much by geometry and topology as by
track fitting alone.

In MicroBooNE the reconstruction proceeds in two passes. The first,
\emph{PandoraCosmic}, targets through-going cosmic-ray activity and removes
hits that are confidently assigned to cosmic topologies. The second,
\emph{PandoraNu}, runs on the residual hit collection and attempts to construct
neutrino-candidate slices, interaction vertices, and outgoing particle
hypotheses. This ordering is important. At the surface, long cosmic tracks are
the dominant source of combinatoric confusion, so neutrino reconstruction is
stable only after that background has been suppressed at the hit level.

The principal output is a hierarchy of \texttt{PFParticle} objects with
explicit parent--daughter relations and associations to the underlying clusters,
space points, vertices, tracks, and showers. That hierarchy is the operational
bridge between raw detector information and topology-based analyses.

Figure~\ref{fig:pandora-edm} makes that event-data model explicit. Pandora does
not jump directly from wire signals to fully formed particle hypotheses:
plane-wise hits are first organised into intermediate geometric objects, and
those objects are then attached to \texttt{PFParticle} instances linked into a
reconstruction hierarchy.

\input{figures/pandora-event-data-model}

A representative labelled event display is shown in
Fig.~\ref{fig:microboone-event-display}. It makes the reconstruction task
concrete: charged tracks and electromagnetic activity overlap within a single
wire--time view, and the particle-flow interpretation must therefore be
assembled from ambiguous projective information rather than read off directly
from the image.

\begin{figure}[htbp!]
  \centering
  \includegraphics[width=0.76\textwidth]{figures/microboone-event-display.png}
  \caption[Example MicroBooNE event display with reconstructed particle labels.]{Example MicroBooNE event display with reconstructed particle labels for charged tracks and electromagnetic activity. The display overlays per-pixel class labels on a calibrated wire--time view from a reconstructed event, highlighting track-like and shower-like signatures. The scene illustrates the complex, overlapping topologies that motivate multi-view reconstruction and image-based interpretation. The display is qualitative and does not encode absolute charge, energy, or timing scales.}
  \label{fig:microboone-event-display}
\end{figure}

It also
defines the limits of an object-level strategy for detached strange-particle
decays: Pandora is optimised for prompt charged-current neutrino interactions,
so displaced \(\Lambda\to p\pi^{-}\) topologies are not guaranteed to appear as
clean canonical particle-flow objects. This limitation is one of the reasons
the present thesis treats detector-level summaries, rather than only
object-level reconstructions, as the primary carriers of information for the
strange-baryon measurement.
